\begin{table}[htbp]
\TABLE{Predeclared residual and cycle-detection protocol (reviewer item C.25). Every value was fixed before any outcome was inspected.\label{tab:fpprotocol}}
{\begin{tabular}{ll}
\toprule
quantity & value \\
\midrule
residual norm & normalized $\ell_1$: $\|T^k_{\theta}(p)-p\|_1/\sum_i|D_i|$ \\
residual threshold & 1e-09 \\
continuation horizon & 200 frozen epochs \\
cycle periods tested & $k \in \{2,3,4\}$ \\
perturbation magnitude & 0.001 (logit coordinates) \\
number of perturbations & 8 \\
return threshold & 1e-06 \\
floating-point precision & IEEE-754 binary64 \\
annealed epoch budget & 300 \\
decode frequency & every 5 epochs \\
\bottomrule
\end{tabular}}
{A state is called an approximate fixed point only when $r_1\le$ the residual threshold. A state is flagged as a period-$k$ candidate only when $r_k$ is below the threshold while $r_d$ remains at least two orders of magnitude larger for every proper divisor $d$ of $k$, including $d=1$; this excludes fixed points, which have small $r_k$ for every $k$, from being reported as short cycles.}
\end{table}
